\documentclass[a4paper]{article}
\usepackage[affil-it]{authblk}
\usepackage[backend=bibtex,style=numeric]{biblatex}
\usepackage{graphicx}
\usepackage{amsmath}
\usepackage{amsfonts}
\usepackage{amssymb}
\usepackage{geometry}
\usepackage{listings}
\usepackage{xcolor}

\lstset{
  language=C++,
  basicstyle=\ttfamily\small,
  keywordstyle=\color{blue},
  commentstyle=\color{gray},
  numberstyle=\tiny,
  numbers=left,
  frame=single,
  breaklines=true
}
\geometry{margin=1.5cm, vmargin={0pt,1cm}}
\setlength{\topmargin}{-1cm}
\setlength{\paperheight}{29.7cm}
\setlength{\textheight}{25.3cm}

\addbibresource{citation.bib}

\begin{document}
% =================================================
\title{Numerical Analysis homework 2}

\author{Han Xiao; 3230100653
  \thanks{Electronic address: \texttt{3230100653@zju.edu.cn}}}
\affil{(Automation 2302), Zhejiang University }


\date{Due time: \today}

\maketitle



% ============================================
\section*{I. }
(Due to the complicated fractions and equations, it's too troublesome to type them in \LaTeX. Therefore, I solve them by hand.)
\begin{figure}[h]
  \centering
  \includegraphics[width=\textwidth]{I.jpg}
  \caption{Problem I}
  \label{fig:I}
\end{figure}


\section*{II. }
Refer to the Lagrange formula, the basis  $l_k(x) = \displaystyle\prod_{i=0, i\neq k}^{n}\cfrac{x-x_i}{x_k-x_i}$.

We've been given $(n+1)$ points to construct a polynomial of degree $2n$, also the polynomial needs to be non-negative in the whole real line.
Therefore, we can use squares of $l_k(x)$ to construct the polynomial:
\[
l_k(x) = \displaystyle\prod_{i=0, i\neq k}^{n}\cfrac{(x-x_i)^2}{(x_k-x_i)^2},
\]
which still satisfies $l_k(x_i)= 0,\quad i \neq k$ and $l_k(x_k) = 1$, while the degree is $2n$.

So the function we want is
\[
p(x) = \sum_{k=0}^{n}f(x_k)l_k(x) = \sum_{k=0}^{n}f(x_k)\displaystyle\prod_{i=0, i\neq k}^{n}\cfrac{(x-x_i)^2}{(x_k-x_i)^2}
\]

\section*{III. }
\begin{figure}[h]
  \centering
  \includegraphics[width=\textwidth]{III-a.jpg}
  \caption{Problem III-a}
  \label{fig:III-a}
\end{figure}

\subsection*{III-b}
For $t = 0$, we have $f[0,1,\dots,n] = \cfrac{(e-1)^n}{n!}$.
We need to find out the $\xi $ in $[0,n]$ such that $f[0,1,\dots,n] = \cfrac{f^{(n)}(\xi)}{n!}$.
Since $f(x) = e^x$, we have $f^{(n)}(x) = e^x$. Therefore, we have $e^{\xi} = (e-1)^n$.
Do \textit{log} on both sides, we have 
\[
\xi = n\ln(e-1)
\].

Roughly, we can calculate that $\ln(e-1) \approx 0.541 > 0.5$.

Therefore, $\xi > \cfrac{1}{2} n$, in other words, located to the\textbf{ right side of the midpoint}.

\section*{IV. }
\subsection*{IV-a}
Constructing the table below:
\begin{figure}[h]
  \centering
  \includegraphics[width=0.3\textwidth]{IV-a.jpg}
  \caption{table of divided differences}
  \label{fig:IV}
\end{figure}
So we have
\[
p_3(f;x) = 5 - 2x + x(x-1) + 0.25x(x-1)(x-2) = 0.25x^3 + 0.25x^2 - 2.5x + 5
\]

\subsection*{IV-b}
We have the approximate derivative of $f(x)$:

\[
f'(x) \approx p_3'(f;x) = 0.75x^2 + 0.5x - 2.5 = 0.25(3x^2 + 2x - 10)
\]
Let $p_3'(f;x) = 0$, we have $x = \cfrac{-1 \pm \sqrt{31}}{3}$. Only $x = \cfrac{-1 + \sqrt{31}}{3} \approx 1.523$ is in the interval $(1,3)$.
\[
\therefore \hat{x}_{min} = 1.523
\]

\section*{V. }
\subsection*{V-a}
For $f(x) = x^7$, we have $f'(x) = 7x^6$, $f''(x) = 42x^5$,
Constructing the table below:
\begin{figure}[h]
  \centering
  \includegraphics[width=0.35\textwidth]{V-a.jpg}
  \caption{table of divided differences}
  \label{fig:V-a}
\end{figure}
\[
\therefore f[0,1,1,1,2,2] = 51
\]

\subsection*{V-b}
According to \textbf{Theorem 2.22}, we have $f[x_0,x_1 \dots x_n] = \cfrac{1}{n!}f^{(n)}(\xi)$.
Here $n = 5$, so we have
\[
f[x_0,x_1,x_2,x_3,x_4,x_5] = f[0,1,1,1,2,2] =
\cfrac{1}{5!}f^{(5)}(\xi) = 51
\]
For $f(x) = x^7$, we have $f^{(5)}(x) = 2520x^2$.
So we have 
\[
\cfrac{1}{5!}f^{(5)}(\xi) = 21{\xi}^2 = 51
\]

\[
\therefore \xi = \sqrt{\cfrac{17}{7}} \approx 1.558 \in (0,2)
\]

\section*{VI. }
\subsection*{VI-a}  
Based on Hermite interpolation rules, constructing the table below:
\begin{figure}
  \centering
  \includegraphics[width=0.4\textwidth]{VI-a.jpg}
  \caption{table of divided differences}
  \label{fig:VI-a}
\end{figure}
So we have
\begin{align*}
f(x) &\approx p(f;x) = 1 + x - 2x(x-1) + \cfrac{2}{3}x(x-1)^2 - \cfrac{5}{36}x(x-1)^2(x-3)\\
&= 1 + x(1+(x-1)(-2 + (x-1)(\cfrac{2}{3} - \cfrac{5}{36}(x-3))))\\
\end{align*}
\[
\therefore f(2) \approx \cfrac{11}{18} \approx 0.6111
\]

\subsection*{VI-b}
According to \textbf{Theorem 2.7}, we have $R_n(f;x) = \cfrac{f^{(n+1)}(\xi)}{(n+1)!}\Pi^n_{i=0}(x-x_i)$.
\[
\therefore R_4(f;x) = \cfrac{f^{(5)}(\xi)}{5!}\Pi^4_{i=0}(x-x_i),\quad \xi \in (0,3)
\]

\[
\because  |f^(5)(x)| \leq M, \quad x \in [0,3]
\]

\[
\therefore |R_4(f;x)| \leq \cfrac{M}{5!}|x(x-1)^2(x-3)|, \quad x \in [0,3]
\]
Now, we have to find out the maximum of $|x(x-1)^2(x-3)|$ in $[0,3]$. The shape of the curve is shown in figure \ref{fig:VI-b}.
\begin{figure}[h]
  \centering
  \includegraphics[width=0.5\textwidth]{VI-b.png}
  \caption{Curve of $y = x(x-1)^2(x-3)$}
  \label{fig:VI-b}
\end{figure}

By observation, we can see that the maximum is roughly located at $x = 2.4$.
For $y = x(x-1)^2(x-3) = x^4 - 5x^3 + 7x^2 -3x$, we have $y' = 4x^3 - 15x^2 + 14x - 3$.
Use \textbf{Newton's method} in the 1st programming homework to find out the root of $y' = 0$:
\begin{lstlisting}
    NewtonSolver(Function &F) : F(F) {}
    double solve(double x0)
    {
        int M = 5;
        double x = x0;
        int k = 0;
        for (; k < M; k++)
        {
            double u = F(x);
            x -= u / (F.d(x));
        }
        return x;
    }

class Func1 : public Function
{
public:
    double operator()(double x)
    {
        return 4*pow(x,3) -15*pow(x,2) + 14*x - 3;
    }
    double d(double x)
    {
        return 12*pow(x,2) - 30*x + 14;
    }
};

int main()
{
    Func1 func;
    NewtonSolver solver(func);
    double x = solver.solve(2.4);
    cout << x << endl;
    return 0;
}
\end{lstlisting}
After running the code, we get $x = 2.443$.
\[
\therefore |x(x-1)^2(x-3)| \leq 2.8334,\quad x \in [0,3]
\]
\begin{align*}
  \therefore |R_4(f;x)| &\leq \cfrac{2.8334M}{120} = 0.02361M, \quad x \in [0,3]
\end{align*}
the maximum possible error 
\begin{align*}
  \epsilon &= 0.02361M
\end{align*}  

\newpage
\section*{VII. }
\begin{figure}[h]
  \centering
  \includegraphics[width=0.9\textwidth]{VII.jpg}
  \caption{Problem VII}
  \label{fig:VII}
\end{figure}

\section*{VIII. }
The definition of partial derivative is:
\[
\cfrac{\partial f}{\partial x_i} = \lim_{h \to 0}\cfrac{f(x_1, x_2, \dots, x_i + h, \dots, x_n) - f(x_1, x_2, \dots, x_i, \dots, x_n)}{h}
\]
According to it, viewing $x_o$ as a variable,
\[
\therefore \cfrac{\partial }{\partial x_0}f[x_0, x_1, \dots, x_n] = \lim_{h \to 0}\cfrac{1}{h}[f(x_0 + h, x_1, \dots, x_n) - f(x_0, x_1, \dots, x_n)]
\]
Based on the definition of divided difference, we have
\[
\cfrac{\partial }{\partial x_0}f[x_0, x_1, \dots, x_n] = f[x_0, x_0, x_1, \dots, x_n]
\]

\section*{IX. }
According to \textbf{Theorem 2.47}, the Chebyshev polynomial has the smallest maximum norm on $[-1,1]$ among polynomials of degree $n$.
namely,
\[
\max_{-1 \leq x \leq 1}|\cfrac{T_n(x)}{2^{n-1}}| \leq \max_{-1 \leq x \leq 1}|p(x)|
\]

Now, we want to find the similar polynomial $a_0x^n + a_1x^{n-1} + \dots + a_n$ on $[a,b]$.
So we need to transform $x \in [a,b]$ to $t \in [-1,1]$: 
\[
t = \cfrac{b-a}{2}x + \cfrac{a+b}{2}
\]
when $x = -1$, we have $t = a$; and when $x = 1$, we have $t = b$.
\[
\therefore x = \cfrac{2t-a-b}{b-a}
\]
Therefore, we have the wanted polynomial on $x \in [a,b]$:
\[
a_0x^n + a_1x^{n-1} + \dots + a_n = a_0 \cfrac{T_n(\cfrac{2x-a-b}{b-a})}{2^{n-1}}
\]
Just sort out the coefficients of $x^n, x^{n-1}, \dots, x^0$ and can simply get $a_0, a_1, \dots, a_n$.

\section*{X. }
We need to prove that
\[
\max_{-1 \leq x \leq 1}|\cfrac{T_n(x)}{T_n(a)}| \leq \max_{-1 \leq x \leq 1}|p(x)|
\]
Proof by contradiction: Suppose that
\[
\exists p \in \mathbb{P}^{a}_n\quad s.t. \quad |\cfrac{T_n(x)}{T_n(a)}| > \max_{-1 \leq x \leq 1}|p(x)|
\]
and 
\[
\because |T_n(x)| \leq 1, \quad x \in [-1,1],
\]
therefore, we only need 
\begin{equation}
  \label{eq:*}
  \exists p \in \mathbb{P}^{a}_n,\quad s.t. \quad |\cfrac{1}{T_n(a)}| > \max_{-1 \leq x \leq 1}|p(x)|.
\end{equation}
Construct a new polynomial $Q(x) = T_n(x) - T_n(a)p(x)$, whose degree is at most $n$.

Since $T_n(x)$ has $n+1$ extremal points in $[-1,1]$, denoted as $x^{'}_0, x^{'}_1, \dots, x^{'}_n$, and at these points $|T_n(x_i)| = 1$.

So we have 
\[
Q(x^{'}_k) = T_n(x^{'}_k) - T_n(a)p(x^{'}_k) = (-1)^k - T_n(a)p(x^{'}_k), \quad k = 0,1,\dots,n.
\]
When the $k$ is even, due to \textbf{inequality \ref{eq:*}} we have $Q(x^{'}_k) = 1 - T_n(a)p(x^{'}_k) > 0$; when the $k$ is odd, we have $Q(x^{'}_k) = -1 - T_n(a)p(x^{'}_k) < 0$.

Since there are $n+1$ extrema, $Q(x)$ has at least $n$ zeros in $[-1,1]$.
While $p \in \mathbb{P}^{a}_n$, so $Q(a) = T_n(a) - T_n(a)p(a) = 1-1 = 0$.
In total, $Q(x)$ has at least $n+1$ zeros in $[-1,1]$, and the degree of $Q(x)$ is at most $n$.
Therefore, $Q(x) \equiv 0$, namely $\cfrac{T_n(x)}{T_n(a)} \equiv p(x)$, which is contradictory to \textbf{inequality \ref{eq:*}}.

\section*{XI.}
The Bernstein base polynomial is defined as 
\[
b^{n}_k(t) = \binom{n}{k}t^k(1-t)^{n-k}, \quad k = 0,1,\dots,n.
\]
\begin{figure}[h]
  \centering
  \includegraphics[width=0.85\textwidth]{XI.jpg}
  \caption{Problem XI}
  \label{fig:XI}
\end{figure}
The Conbinatorial number $\binom{n}{k}$ has the property that $\binom{n}{k} = \binom{n-1}{k-1} + \binom{n-1}{k}$.
Based on this property, we have: 

\section*{XII.}
The Bernstein base polynomial is defined as 
\[
b^{n}_k(t) = \binom{n}{k}t^k(1-t)^{n-k}, \quad k = 0,1,\dots,n.
\]
\[
\therefore \int_{0}^{1} b^{n}_k(t) \,dt = \binom{n}{k} \int_{0}^{1} t^k(1-t)^{n-k} \,dt
\]
According to the definition of Beta function, we have
\[
B(p,q) = \int_{0}^{1} t^{p-1}(1-t)^{q-1} \,dt
\]
Let $p = k+1, \quad q = n-k+1$. And Beta function has the property that
\[
B(p,q) = \cfrac{\Gamma(p)\Gamma(q)}{\Gamma(p+q)}
\]
For $k \in \mathbb{Z}, \quad \Gamma(k) = (k-1)!$.
\[
\therefore \int_{0}^{1} b^{n}_k(t) \,dt = \binom{n}{k} \cfrac{\Gamma(k+1)\Gamma(n-k+1)}{\Gamma(n+2)} = \binom{n}{k} \cfrac{k!(n-k)!}{(n+1)!} = \cfrac{1}{n+1}
\]

\end{document}